% Options for packages loaded elsewhere
\PassOptionsToPackage{unicode}{hyperref}
\PassOptionsToPackage{hyphens}{url}
\documentclass[
]{article}
\usepackage{xcolor}
\usepackage[margin=1in]{geometry}
\usepackage{amsmath,amssymb}
\setcounter{secnumdepth}{-\maxdimen} % remove section numbering
\usepackage{iftex}
\ifPDFTeX
  \usepackage[T1]{fontenc}
  \usepackage[utf8]{inputenc}
  \usepackage{textcomp} % provide euro and other symbols
\else % if luatex or xetex
  \usepackage{unicode-math} % this also loads fontspec
  \defaultfontfeatures{Scale=MatchLowercase}
  \defaultfontfeatures[\rmfamily]{Ligatures=TeX,Scale=1}
\fi
\usepackage{lmodern}
\ifPDFTeX\else
  % xetex/luatex font selection
\fi
% Use upquote if available, for straight quotes in verbatim environments
\IfFileExists{upquote.sty}{\usepackage{upquote}}{}
\IfFileExists{microtype.sty}{% use microtype if available
  \usepackage[]{microtype}
  \UseMicrotypeSet[protrusion]{basicmath} % disable protrusion for tt fonts
}{}
\makeatletter
\@ifundefined{KOMAClassName}{% if non-KOMA class
  \IfFileExists{parskip.sty}{%
    \usepackage{parskip}
  }{% else
    \setlength{\parindent}{0pt}
    \setlength{\parskip}{6pt plus 2pt minus 1pt}}
}{% if KOMA class
  \KOMAoptions{parskip=half}}
\makeatother
\usepackage{color}
\usepackage{fancyvrb}
\newcommand{\VerbBar}{|}
\newcommand{\VERB}{\Verb[commandchars=\\\{\}]}
\DefineVerbatimEnvironment{Highlighting}{Verbatim}{commandchars=\\\{\}}
% Add ',fontsize=\small' for more characters per line
\usepackage{framed}
\definecolor{shadecolor}{RGB}{248,248,248}
\newenvironment{Shaded}{\begin{snugshade}}{\end{snugshade}}
\newcommand{\AlertTok}[1]{\textcolor[rgb]{0.94,0.16,0.16}{#1}}
\newcommand{\AnnotationTok}[1]{\textcolor[rgb]{0.56,0.35,0.01}{\textbf{\textit{#1}}}}
\newcommand{\AttributeTok}[1]{\textcolor[rgb]{0.13,0.29,0.53}{#1}}
\newcommand{\BaseNTok}[1]{\textcolor[rgb]{0.00,0.00,0.81}{#1}}
\newcommand{\BuiltInTok}[1]{#1}
\newcommand{\CharTok}[1]{\textcolor[rgb]{0.31,0.60,0.02}{#1}}
\newcommand{\CommentTok}[1]{\textcolor[rgb]{0.56,0.35,0.01}{\textit{#1}}}
\newcommand{\CommentVarTok}[1]{\textcolor[rgb]{0.56,0.35,0.01}{\textbf{\textit{#1}}}}
\newcommand{\ConstantTok}[1]{\textcolor[rgb]{0.56,0.35,0.01}{#1}}
\newcommand{\ControlFlowTok}[1]{\textcolor[rgb]{0.13,0.29,0.53}{\textbf{#1}}}
\newcommand{\DataTypeTok}[1]{\textcolor[rgb]{0.13,0.29,0.53}{#1}}
\newcommand{\DecValTok}[1]{\textcolor[rgb]{0.00,0.00,0.81}{#1}}
\newcommand{\DocumentationTok}[1]{\textcolor[rgb]{0.56,0.35,0.01}{\textbf{\textit{#1}}}}
\newcommand{\ErrorTok}[1]{\textcolor[rgb]{0.64,0.00,0.00}{\textbf{#1}}}
\newcommand{\ExtensionTok}[1]{#1}
\newcommand{\FloatTok}[1]{\textcolor[rgb]{0.00,0.00,0.81}{#1}}
\newcommand{\FunctionTok}[1]{\textcolor[rgb]{0.13,0.29,0.53}{\textbf{#1}}}
\newcommand{\ImportTok}[1]{#1}
\newcommand{\InformationTok}[1]{\textcolor[rgb]{0.56,0.35,0.01}{\textbf{\textit{#1}}}}
\newcommand{\KeywordTok}[1]{\textcolor[rgb]{0.13,0.29,0.53}{\textbf{#1}}}
\newcommand{\NormalTok}[1]{#1}
\newcommand{\OperatorTok}[1]{\textcolor[rgb]{0.81,0.36,0.00}{\textbf{#1}}}
\newcommand{\OtherTok}[1]{\textcolor[rgb]{0.56,0.35,0.01}{#1}}
\newcommand{\PreprocessorTok}[1]{\textcolor[rgb]{0.56,0.35,0.01}{\textit{#1}}}
\newcommand{\RegionMarkerTok}[1]{#1}
\newcommand{\SpecialCharTok}[1]{\textcolor[rgb]{0.81,0.36,0.00}{\textbf{#1}}}
\newcommand{\SpecialStringTok}[1]{\textcolor[rgb]{0.31,0.60,0.02}{#1}}
\newcommand{\StringTok}[1]{\textcolor[rgb]{0.31,0.60,0.02}{#1}}
\newcommand{\VariableTok}[1]{\textcolor[rgb]{0.00,0.00,0.00}{#1}}
\newcommand{\VerbatimStringTok}[1]{\textcolor[rgb]{0.31,0.60,0.02}{#1}}
\newcommand{\WarningTok}[1]{\textcolor[rgb]{0.56,0.35,0.01}{\textbf{\textit{#1}}}}
\usepackage{graphicx}
\makeatletter
\newsavebox\pandoc@box
\newcommand*\pandocbounded[1]{% scales image to fit in text height/width
  \sbox\pandoc@box{#1}%
  \Gscale@div\@tempa{\textheight}{\dimexpr\ht\pandoc@box+\dp\pandoc@box\relax}%
  \Gscale@div\@tempb{\linewidth}{\wd\pandoc@box}%
  \ifdim\@tempb\p@<\@tempa\p@\let\@tempa\@tempb\fi% select the smaller of both
  \ifdim\@tempa\p@<\p@\scalebox{\@tempa}{\usebox\pandoc@box}%
  \else\usebox{\pandoc@box}%
  \fi%
}
% Set default figure placement to htbp
\def\fps@figure{htbp}
\makeatother
\setlength{\emergencystretch}{3em} % prevent overfull lines
\providecommand{\tightlist}{%
  \setlength{\itemsep}{0pt}\setlength{\parskip}{0pt}}
\usepackage{bookmark}
\IfFileExists{xurl.sty}{\usepackage{xurl}}{} % add URL line breaks if available
\urlstyle{same}
\hypersetup{
  pdftitle={Bayesian Mixed Models},
  hidelinks,
  pdfcreator={LaTeX via pandoc}}

\title{Bayesian Mixed Models}
\author{}
\date{\vspace{-2.5em}2026-04-17}

\begin{document}
\maketitle

\subsection{Introduction}\label{introduction}

Bayesian mixed models treat the random-effect variance components as
parameters with their own priors instead of plug-in estimates from REML.
The pay-off is most visible in small samples and at the boundary: ML and
REML routinely return zero variance estimates when only a handful of
clusters contribute, but a half-Normal or half-Cauchy prior keeps the
variance away from zero in a principled way and propagates the resulting
uncertainty into every fixed-effect interval. The cost is computational
--- sampling instead of optimising --- but for any analysis that depends
on credible variance components, that cost is a worthwhile investment.

\subsection{Prerequisites}\label{prerequisites}

A working understanding of frequentist mixed-effects models
(\texttt{lme4}-style formulas), Bayesian regression, and the role of
half-distributions as priors on standard deviations.

\subsection{Theory}\label{theory}

A random-intercept-and-slope model has the form

\[y_{ij} = \beta_0 + \beta_1 X_{ij} + u_{0i} + u_{1i} X_{ij} + \varepsilon_{ij},\]

with \(\bigl(u_{0i}, u_{1i}\bigr) \sim \mathcal N(0, \Sigma)\) and
\(\varepsilon_{ij} \sim \mathcal N(0, \sigma^2)\). Priors are placed on
every component: weakly informative Normal priors on \(\beta\),
half-Normal or half-Cauchy on the diagonal of \(\Sigma\), an LKJ prior
on the correlation, and an exponential or half-\(t\) on \(\sigma\). The
LKJ\((\eta)\) prior with \(\eta > 1\) pulls correlations toward zero
gently while still allowing strong correlations when the data demand
them. Posterior draws of \(u_i\) are \emph{shrunk} toward zero by an
amount that depends on within-cluster signal-to-noise --- the same
partial-pooling behaviour that makes mixed models attractive in the
first place.

\subsection{Assumptions}\label{assumptions}

Conditional independence of observations given random effects,
exchangeability of clusters, approximately Normal random effects with
the prior structure above, and proper priors on every variance and
correlation parameter. Convergence diagnostics (\(\hat R\), divergent
transitions, ESS) are part of the assumptions in practice.

\subsection{R Implementation}\label{r-implementation}

\begin{Shaded}
\begin{Highlighting}[]
\FunctionTok{library}\NormalTok{(brms)}

\FunctionTok{data}\NormalTok{(sleepstudy, }\AttributeTok{package =} \StringTok{"lme4"}\NormalTok{)}

\NormalTok{fit }\OtherTok{\textless{}{-}} \FunctionTok{brm}\NormalTok{(Reaction }\SpecialCharTok{\textasciitilde{}}\NormalTok{ Days }\SpecialCharTok{+}\NormalTok{ (Days }\SpecialCharTok{|}\NormalTok{ Subject), }\AttributeTok{data =}\NormalTok{ sleepstudy,}
           \AttributeTok{prior =} \FunctionTok{prior}\NormalTok{(}\FunctionTok{normal}\NormalTok{(}\DecValTok{0}\NormalTok{, }\DecValTok{30}\NormalTok{), }\AttributeTok{class =} \StringTok{"b"}\NormalTok{) }\SpecialCharTok{+}
                   \FunctionTok{prior}\NormalTok{(}\FunctionTok{normal}\NormalTok{(}\DecValTok{250}\NormalTok{, }\DecValTok{50}\NormalTok{), }\AttributeTok{class =} \StringTok{"Intercept"}\NormalTok{) }\SpecialCharTok{+}
                   \FunctionTok{prior}\NormalTok{(}\FunctionTok{exponential}\NormalTok{(}\FloatTok{0.1}\NormalTok{), }\AttributeTok{class =} \StringTok{"sd"}\NormalTok{) }\SpecialCharTok{+}
                   \FunctionTok{prior}\NormalTok{(}\FunctionTok{exponential}\NormalTok{(}\FloatTok{0.1}\NormalTok{), }\AttributeTok{class =} \StringTok{"sigma"}\NormalTok{) }\SpecialCharTok{+}
                   \FunctionTok{prior}\NormalTok{(}\FunctionTok{lkj}\NormalTok{(}\DecValTok{2}\NormalTok{), }\AttributeTok{class =} \StringTok{"cor"}\NormalTok{),}
           \AttributeTok{chains =} \DecValTok{4}\NormalTok{, }\AttributeTok{iter =} \DecValTok{2000}\NormalTok{, }\AttributeTok{refresh =} \DecValTok{0}\NormalTok{)}
\FunctionTok{summary}\NormalTok{(fit)}

\CommentTok{\# Visualise random effects}
\FunctionTok{ranef}\NormalTok{(fit)}
\end{Highlighting}
\end{Shaded}

\subsection{Output \& Results}\label{output-results}

\texttt{summary(fit)} reports the population-level fixed effects, the
standard deviations of the random intercept and slope, the correlation
between them, and the residual scale, all with credible intervals,
\(\hat R\), and ESS. \texttt{ranef(fit)} returns a list of per-cluster
posterior summaries --- the analogue of \texttt{lme4::ranef()} but with
full credible intervals.

\subsection{Interpretation}\label{interpretation}

A reporting sentence: ``A Bayesian linear mixed model estimated the
population-level slope at 10.5 ms per day of sleep deprivation (95 \%
credible interval 7.4 to 13.6); between-subject heterogeneity in slope
was substantial (\(\mathrm{SD} = 6.0\), 95 \% CrI 4.0 to 8.8) and weakly
negatively correlated with the intercept (\(r = -0.13\), 95 \% CrI
\(-0.49\) to 0.27).'' Always report the variance components with their
credible intervals; point estimates alone can be misleading when
posteriors are skewed.

\subsection{Practical Tips}\label{practical-tips}

\begin{itemize}
\tightlist
\item
  Use \texttt{LKJ(2)} for correlation priors as a sensible default;
  tighten to \texttt{LKJ(4)} only when you have strong reason to expect
  weak correlations.
\item
  For divergent transitions in hierarchical models, switch to
  non-centred parameterisation (\texttt{brms} does this automatically
  when you set \texttt{prior(...\ ,\ class\ =\ "sd",\ ...)}).
\item
  Standardise continuous predictors before fitting; weakly informative
  priors are then meaningful across columns.
\item
  Compare nested random-effect structures with \texttt{loo\_compare()}
  rather than likelihood-ratio tests; LOO works correctly when ML-LRT
  statistics are on the boundary of the parameter space.
\item
  For very small \(n_{\text{cluster}}\) (under 5), report a sensitivity
  analysis across at least two priors on the random-effect SD; the
  choice can dominate the posterior.
\item
  \texttt{rstanarm::stan\_lmer()} mirrors \texttt{lme4} syntax and uses
  pre-compiled Stan models, useful when fit time matters and the model
  is standard.
\end{itemize}

\subsection{R Packages Used}\label{r-packages-used}

\texttt{brms} for formula-based Bayesian mixed models with full prior
control, \texttt{rstanarm} for the pre-compiled \texttt{stan\_lmer}
interface, \texttt{bayesplot} for diagnostic plots, and
\texttt{tidybayes} for tidy posterior draws when bespoke shrinkage plots
are needed.

\subsection{For Reviewers}\label{for-reviewers}

\textbf{What to look for in a paper using this method.}

\begin{itemize}
\tightlist
\item
  \textbf{Common misapplications.}
\item
  \textbf{Diagnostics that should be reported but often aren't.}
\item
  \textbf{Red flags in tables and figures.}
\item
  \textbf{What to verify.}
\item
  \textbf{What an adequate Methods paragraph must contain.}
\end{itemize}

\subsection{See also --- labs in this
chapter}\label{see-also-labs-in-this-chapter}

\begin{itemize}
\tightlist
\item
  \href{labs/lab-bayesian-thinking-control-vs-decision-error.Rmd}{Bayesian
  thinking; α control vs decision error}
\item
  \href{labs/lab-brms-stan-loo-hierarchical-models.Rmd}{brms / Stan,
  LOO, hierarchical models}
\end{itemize}

\end{document}
